% !TEX program = xelatex
% ESP32 车载性能仪 —— 个人技术知识库
% 编译：xelatex KNOWLEDGE_BASE.tex （需运行两次以生成目录/书签）
% 依赖等宽字体 Consolas（Windows 自带，含制表符字形，用于 ASCII 架构图）
\documentclass[UTF8,11pt,a4paper]{ctexart}

\usepackage[top=2cm,bottom=2cm,left=2cm,right=2cm]{geometry}
\usepackage{amsmath,amssymb}
\usepackage{xcolor}
\usepackage{enumitem}
\usepackage{booktabs}
\usepackage{tabularx,array,longtable}
\usepackage{fancyvrb}
\usepackage{tcolorbox}
\tcbuselibrary{skins,breakable}
\usepackage{underscore} % 让正文中的 _ 直接显示，省去大量转义
\usepackage{fancyhdr}
\usepackage{titlesec}
\usepackage[colorlinks=true,linkcolor=teal!70!black,urlcolor=blue!60!black,
            bookmarksnumbered=true]{hyperref}

% 等宽字体（含 Unicode 制表符，用于 ASCII 图）
\setmonofont{Consolas}

% 页眉页脚
\pagestyle{fancy}\fancyhf{}
\fancyhead[L]{\small\color{gray} ESP32 车载性能仪 · 技术知识库}
\fancyhead[R]{\small\color{gray}\leftmark}
\fancyfoot[C]{\small\thepage}
\renewcommand{\headrulewidth}{0.3pt}

% 章节样式
\titleformat{\section}{\Large\bfseries\color{teal!70!black}}
  {\thesection}{0.6em}{}[{\color{teal!70!black}\titlerule[1.2pt]}]
\titlespacing{\section}{0pt}{14pt}{6pt}
\titleformat{\subsection}{\large\bfseries\color{teal!60!black}}{\thesubsection}{0.5em}{}
\titlespacing{\subsection}{0pt}{8pt}{3pt}
\titleformat{\subsubsection}{\normalsize\bfseries\color{teal!55!black}}{}{0pt}{}
\titlespacing{\subsubsection}{0pt}{6pt}{2pt}

\setlist{topsep=2pt,itemsep=1pt,leftmargin=1.6em}

% 强调框 / 警告框
\newtcolorbox{say}[1]{enhanced,breakable,
  colback=teal!5,colframe=teal!60!black,arc=1.2mm,boxrule=0.6pt,
  left=6pt,right=6pt,top=4pt,bottom=4pt,
  fonttitle=\bfseries\color{white},colbacktitle=teal!60!black,title={#1}}
\newtcolorbox{warn}[1]{enhanced,breakable,
  colback=red!4,colframe=red!55!black,arc=1.2mm,boxrule=0.6pt,
  left=6pt,right=6pt,top=4pt,bottom=4pt,
  fonttitle=\bfseries\color{white},colbacktitle=red!55!black,title={#1}}

% ASCII 架构图环境（等宽 + 浅色边框）
\DefineVerbatimEnvironment{ascii}{Verbatim}%
  {fontsize=\footnotesize,frame=single,framesep=4pt,rulecolor=\color{teal!40},
   commandchars=\\\{\}}

\renewcommand{\arraystretch}{1.25}

\title{\textbf{\Huge Embedded Knowledge Base}\\[6pt]
       \Large ESP32 车载性能仪 · 个人技术知识库}
\author{工作区：\texttt{e:\textbackslash esp32\textbackslash esp32-imu}}
\date{生成日期：2026-06-10}

\begin{document}
\maketitle

\begin{say}{这份文档是写给半年后的自己的}
重新打开这个仓库时，\textbf{只读这一份文档}就能想起来：当时做了什么、为什么这么设计、
数据怎么流、核心技术是什么、踩过哪些坑。
\end{say}

\vspace{4pt}
\noindent\textbf{关于工作区范围：}当前工作区只有\textbf{一个固件项目}（\texttt{esp32-imu}），
外加随附文档（GPS 升级分析、面试归档、接线图、构建指南）和 \texttt{cad/} 里的 OpenSCAD 外壳模型
与 Gerber 压缩包。本文档据此组织。

\tableofcontents
\newpage

% ================================================================
\section{项目全景图}
% ================================================================

\subsection{项目名称：ESP32 车载性能仪（“Dragy 风格” Performance Box）}

\subsubsection{项目背景}
想用便宜的开源硬件，自己做一台\textbf{摩托车/汽车的车载性能记录仪}，对标商业产品 Dragy
（一种 GPS 性能盒子，专测 0--100 km/h、0--400m、刹车距离）。商业 Dragy 要 2000 元，
而一套 ESP32 + MPU6050 + GT-U7 GPS 的料才几十块。

核心挑战不是“读个传感器显示出来”，而是：
\begin{itemize}
  \item \textbf{传感器装在摩托车上，发动机振动极烈}，MPU6050 原始数据被 20\textasciitilde 50 Hz 机械振动淹没。
  \item \textbf{GT-U7 是 30 元级老 GPS 模块}，静止零漂 ±2 km/h，10 Hz 更新，低速段噪声比信号还大。
  \item 要在这种“脏数据”上算出能看的 G 力、倾角、加速测试成绩。
\end{itemize}
项目真正的价值在于\textbf{用软件算法把廉价硬件的潜力榨干}。

\subsubsection{项目目标}
\begin{enumerate}
  \item 实时显示\textbf{速度 / G 力（纵向、横向、垂直）/ 压弯倾角 / 加速测试成绩}。
  \item 自动识别“开始骑行 / 停车”，自动开始/停止数据记录与录音。
  \item 通过 \textbf{WiFi 热点 + 网页}把数据推到手机浏览器（无需 App）。
  \item 专门的\textbf{摩托车模式}：压弯角度、翘头/翘尾检测、弯道半径估算。
  \item 数据落盘（SPIFFS 录音 / SD 卡 CSV 行车数据）。
\end{enumerate}

\subsubsection{技术栈}
\begin{center}\small
\begin{tabularx}{\linewidth}{@{}lX@{}}
\toprule
\textbf{维度} & \textbf{选型} \\
\midrule
芯片 & ESP32-WROOM-32（双核 240MHz，板载 WiFi） \\
框架 & Arduino Framework（基于 ESP-IDF） \\
操作系统 & 无 RTOS 显式使用，裸机 \texttt{loop()} + 软件分时调度（底层是 FreeRTOS，但未直接用任务/队列） \\
构建系统 & PlatformIO（\texttt{platformio.ini}） \\
语言 & C++（Arduino 风格）+ 前端 HTML/CSS/JS（\texttt{data/index.html}） \\
传感器 & MPU6050（6 轴 IMU）、GT-U7（GPS，UBX-M8030 芯片）、模拟麦克风 \\
显示 & SSD1306 OLED 128×64（I²C） \\
协议 & I²C（双路）、UART（GPS）、SPI（SD 卡 HSPI）、HTTP（Web）、NMEA/PMTK \\
关键库 & Adafruit MPU6050 / SSD1306 / GFX、TinyGPSPlus、ArduinoJson \\
机械 & OpenSCAD 参数化外壳（\texttt{cad/}） \\
\bottomrule
\end{tabularx}
\end{center}

\subsubsection{项目一句话总结（适合放简历）}
\begin{say}{简历版}
基于 ESP32 + MPU6050 + GT-U7 GPS 自研车载性能仪，通过\textbf{三级级联滤波}抑制发动机振动、
\textbf{GPS-IMU 互补滤波融合}与\textbf{动态信心度加权}算法，将 30 元 GPS 模块的静止零漂
从 ±2.5 km/h 压到 ±0.8 km/h；实现 0--100 km/h 加速测试、压弯倾角/翘头检测，
并通过 WiFi 热点 + 网页实时可视化。
\end{say}

\subsubsection{项目目录结构分析}
\begin{ascii}
esp32-imu/
├── platformio.ini              # PlatformIO 工程配置：板型、库依赖、串口
├── src/                        # 固件源码（核心）
│   ├── main.cpp                # 主程序：setup/loop、Web服务器、OLED显示、调度
│   ├── mpu6050.{h,cpp}         # IMU驱动 + 三级振动滤波 + 姿态角计算
│   ├── gps_module.{h,cpp}      # GPS解析 + 零漂抑制 + GPS-IMU融合 + 性能测试
│   ├── performance_analyzer.*  # G力计算 + 摩托车模式（倾角/弯道/翘头）
│   ├── audio_recorder.{h,cpp}  # ADC麦克风录音 → WAV，自动骑行触发
│   └── sd_logger.{h,cpp}       # SD卡 CSV 行车数据记录（HSPI）
├── data/index.html             # SPIFFS 网页前端（手机端仪表盘）
├── cad/                        # 机械外壳（OpenSCAD 参数化模型 + STL）
├── interview/                  # LaTeX 版面试材料（main.tex/pdf）
├── INTERVIEW_NOTES.md          # [*] 面试技术答辩归档：踩坑/原理/量化效果（精华）
├── GPS_UPGRADE_ANALYSIS.md     # GPS 硬件升级对比分析
├── GT-U7_ULTIMATE_OPTIMIZATION.md  # GT-U7 软件压榨优化记录
├── wiring_diagram.txt          # 接线图
├── BUILD_GUIDE.md              # 编译/上传/串口/排错指南
└── README.md                   # 原始 README（[!] 已与代码脱节，见踩坑记录）
\end{ascii}
\noindent\textbf{重点：}真正的“知识密度”集中在 \texttt{src/} 和 \texttt{INTERVIEW\_NOTES.md}。

% ================================================================
\section{系统架构设计}
% ================================================================

\subsection{系统由哪些模块组成}
固件按“传感器采集 → 数据融合 → 输出”分层，五个功能模块围绕 \texttt{main.cpp} 的调度器协作：

\begin{ascii}
┌──────────────────────────────────────────────────────────────┐
│                       main.cpp (调度器)                        │
│   软件分时循环 loop()：按 millis() 时间片驱动各模块            │
└──────────────────────────────────────────────────────────────┘
      │           │            │           │           │
  ┌───▼───┐  ┌────▼───┐   ┌────▼────┐  ┌───▼───┐  ┌────▼───┐
  │mpu6050│  │  gps   │   │perf_ana │  │ audio │  │   sd   │
  │IMU驱动 │  │GPS解析  │   │融合/摩托 │  │ 录音  │  │  记录  │
  │+三级滤波│ │+零漂抑制 │  │G力/倾角 │  │ WAV  │  │  CSV   │
  └───────┘  └────────┘   └─────────┘  └───────┘  └────────┘
      │           │            │
      └───────────┴────────────┘
                  │ 全局数据结构 (mpu6050_data / gps_data / fused_data / moto_data)
                  ▼
        ┌────────────────────────────────┐
        │ 输出层：OLED显示 + 串口 + Web服务器 │
        └────────────────────────────────┘
\end{ascii}

\subsection{为什么这样拆分模块}
\begin{itemize}
  \item \textbf{按物理传感器 + 按算法职责拆分}。\texttt{mpu6050}/\texttt{gps}/\texttt{audio}/\texttt{sd}
        对应四个外设；\texttt{performance\_analyzer} 是纯算法层（不碰硬件，只消费别人的全局数据），
        把“业务逻辑”和“驱动”分开。
  \item \textbf{模块间用全局数据结构通信}。\texttt{mpu6050\_data}、\texttt{gps\_data}、\texttt{fused\_data}、
        \texttt{moto\_data} 都是 \texttt{extern} 全局单例。这是嵌入式裸机常见做法：避免动态分配、零拷贝、
        任何模块随时能读最新值。代价是耦合度高、不可重入——但单线程 \texttt{loop()} 里没有竞态。
  \item \textbf{没有用 FreeRTOS 任务}。所有“并发”靠 \texttt{loop()} 里基于 \texttt{millis()} 的时间片调度
        实现。对这种数据流明确、无阻塞操作的场景，裸机分时比多任务更简单、更可控。
\end{itemize}

\subsection{模块之间如何协作（数据链路）}
\begin{ascii}
MPU6050 (I²C0, 200Hz) ──原始 a/g/temp──► [三级滤波] ──► mpu6050_data.*_Filtered
                                                          │
                                  ┌───────────────────────┤
                                  ▼                        ▼
              fuseIMUSpeed() ◄── GPS速度          [performance_analyzer]
              (互补滤波)                           calculateGForces() ◄ GPS速度差分(信心度)
                  │                                    │
                  ▼                                    ├─► fused_data (G力/状态/功率)
          gps_data.fused_speed                         └─► moto_data  (倾角/弯道/翘头)

GT-U7 GPS (UART2,10Hz) ─NMEA─► TinyGPS++ ─► [零漂抑制 三层] ─► gps_data.speed_kmh / 位置 / 卫星
                                                              │
                                                              ▼  [性能测试状态机] 0-60/0-100/制动

输出汇聚：gps_data + fused_data + moto_data + audio_data
   ├─► OLED (5种显示模式, 20Hz)      ├─► Web /data /moto /audio (JSON)
   └─► 串口 (简要/详细/混合)          └─► SD卡 CSV (10Hz, 18字段/行)
\end{ascii}

\begin{center}\small
\begin{tabularx}{\linewidth}{@{}llX@{}}
\toprule
\textbf{链路} & \textbf{频率} & \textbf{作用} \\
\midrule
IMU 采样 → 滤波 & 200 Hz & 高频采集，为 GPS-IMU 融合提供积分基础 \\
IMU → 速度融合 & 200 Hz & 用加速度积分填补 GPS 帧间空白，GPS 帧到来时拉回漂移 \\
GPS → 零漂抑制 & 10 Hz & 把脏速度变成能用的速度，喂给所有下游 \\
融合 → G力/倾角 & 20 Hz & 算出展示用的 G 力、压弯角、弯道信息 \\
汇聚 → 输出 & 20/10 Hz & OLED/网页/串口（20Hz）、落盘（10Hz） \\
\bottomrule
\end{tabularx}
\end{center}

% ================================================================
\section{项目运行流程}
% ================================================================

\subsection{上电之后发生什么（初始化流程）}
\texttt{setup()}（\texttt{main.cpp:713}）的顺序是踩坑调整过的：
\begin{ascii}
上电/Reset
  → Serial.begin(115200)          串口先起，方便看后续日志
  → pinMode(0, INPUT_PULLUP)      BOOT按钮(GPIO0)做显示模式切换
  → 双路I²C初始化  I2C_MPU(21,22,400k) / I2C_OLED(18,19,400k)
  → scanI2CDevices()              扫总线确认 0x68/0x3C 在位（排错关键）
  → Init_mpu6050()                ±8G/±500°/s，DLPF=94Hz，初始化滤波缓冲
  → initGPS()                     UART2 9600，PMTK配 10Hz+精简NMEA+SBAS
  → initPerformanceAnalyzer() / initMotorcycleMode() / initAudioRecorder()
  → [SD卡初始化被注释跳过]          [!] 当前为“测试模式”（见踩坑记录）
  → display.begin()               OLED启动，开机画面2秒
  → SPIFFS.begin(true)            挂载文件系统（网页+录音）
  → WiFi.softAP("ESP32-IMU")      开热点，固定IP 192.168.4.1
  → server.on(...) + server.begin()   注册6个HTTP路由
  → display_mode=1, displayQRCode()   初始显示WiFi连接信息
\end{ascii}

\subsection{主循环流程（任务调度）}
\texttt{loop()}（\texttt{main.cpp:835}）是系统心跳，用 \texttt{millis()} 时间片做\textbf{软件多任务}：
\begin{ascii}
loop() {
  server.handleClient();        ① Web请求：每圈都处理（非阻塞）
  检测BOOT按钮 → 切换 display_mode (0~4)
  handleSerialCommands();       ② 串口命令：每圈都查
  if (now-lastIMUTime   >= 5ms)   → ReadMPU6050() + fuseIMUSpeed()   // 200Hz
  if (now-lastGPSTime   >= 100ms) → updateGPSData() + SD记录          // 10Hz
  updateAudioRecorder();          → processSample() + 自动录音检查     // 持续
  if (now-lastUpdateTime>= 50ms)  → updateFusedData() + OLED + 串口   // 20Hz
}
\end{ascii}
\begin{say}{关键设计：全程无 delay()}
\texttt{loop()} 里\textbf{绝对没有 delay()}（除了按钮防抖那一处 200ms）。所有周期都用
“时间戳差值比较”实现非阻塞，这样高频 IMU 不会被低频 GPS 拖慢——这是裸机分时调度能跑起来的前提。
\end{say}

\subsection{中断流程}
全项目\textbf{只有一处硬件中断}：录音采样定时器。教科书式的 \textbf{ISR 最小化}。
\begin{ascii}
startRecording()
  → timerBegin(0, 80, true)        定时器0，80分频 → 1MHz计数
  → timerAlarmWrite(timer, 125)    125 tick = 125µs = 8kHz
  → timerAttachInterrupt(onSampleTimer)
        │ 每125µs触发
        ▼
  IRAM_ATTR onSampleTimer() {      中断服务程序（放IRAM加速）
     lastSample = adc1_get_raw()   只做最少的事：读ADC + 置标志位
     sampleReady = true
  }
        │ loop()中
        ▼
  processSample() {                重活放主循环（ISR外）
     8位量化 + 写缓冲 + 满512字节写SPIFFS
  }
\end{ascii}
原因——SPIFFS 写操作可能阻塞，绝不能在 ISR 里做。

\subsection{状态机流程}
项目里有\textbf{三个独立状态机}：

\textbf{1）录音状态机}（\texttt{audio\_recorder.cpp:296}）
\begin{ascii}
RECORDING_IDLE/STOPPED ──速度≥5km/h──► RECORDING_ACTIVE
        ▲                                   │
        │                          速度<2km/h 持续5秒
        └───────────────────────────────────┘
   (也可手动 start/stop/pause/resume)
\end{ascii}

\textbf{2）加速/制动测试状态机}（\texttt{gps\_module.cpp:337}）
\begin{ascii}
速度跨过 5km/h ──► 启动加速测试计时
   ├─ 跨 60km/h  → 记 0-60 时间
   ├─ 跨 100km/h → 记 0-100 时间
   └─ 跨 200km/h → 记 100-200，结束
速度≥95km/h ──► 启动制动测试 ──速度≤5km/h──► 记 100-0 时间
\end{ascii}

\textbf{3）SD 自动记录状态机}（\texttt{sd\_logger.cpp:280}）：速度 $>$5 km/h 开始记录，停车 30 秒后结束。

% ================================================================
\section{核心模块解析}
% ================================================================

\subsection{模块一：MPU6050 IMU 驱动（mpu6050.cpp）}
\textbf{职责}：读 6 轴数据，做振动滤波，算姿态角。\textbf{设计思路}：用 Adafruit 库读原始值，
但\textbf{不信任原始值}——叠加三级软件滤波。\textbf{数据结构} \texttt{MPU6050\_DATA} 同时保存“原始值”
和“\_Filtered 滤波值”两套：下游算 G 力用滤波值，算融合用原始值（积分对延迟敏感）。
\begin{center}\small
\begin{tabularx}{\linewidth}{@{}lX@{}}
\toprule
\textbf{函数} & \textbf{作用} \\
\midrule
\texttt{Init\_mpu6050()} & 量程 ±8G/±500°/s，DLPF=94Hz；失败时自动尝试 0x69 地址 \\
\texttt{ReadMPU6050()} & 读 a/g/temp → 滤波 → 算姿态角 \\
\texttt{FilterDataForVibration()} & 滑动平均(10点) + 一阶 IIR(\(\alpha\)=0.1) 级联 \\
\texttt{CalculateAttitude()} & atan2 算 Roll/Pitch（[!] Yaw 硬编码为 0） \\
\bottomrule
\end{tabularx}
\end{center}

\subsection{模块二：GPS 模块（gps_module.cpp）}
\textbf{职责}：解析 NMEA、抑制零漂、GPS-IMU 速度融合、跑加速/制动测试状态机。
这是\textbf{代码量最大、算法最密集的模块（756 行）}。本质是“清洗 + 融合”流水线。
\begin{center}\small
\begin{tabularx}{\linewidth}{@{}lX@{}}
\toprule
\textbf{函数} & \textbf{作用} \\
\midrule
\texttt{initGPS()} & UART2 + PMTK 配置（10Hz/精简NMEA/SBAS/WAAS） \\
\texttt{updateGPSData()} & 逐字节喂 TinyGPS++ → 三层零漂抑制 → 信号监控自恢复 \\
\texttt{fuseIMUSpeed()} & GPS-IMU 互补滤波（核心算法） \\
\texttt{calibrateIMUBias()} & 静止时递推均值在线校准加速度偏置（O(1) 内存） \\
\texttt{interpolateSpeedAtTime()} & 速度阈值穿越线性插值，把 10Hz 的 100ms 不确定度降到 ~10-20ms \\
\texttt{calculateDistance()} & Haversine 球面距离（位置跳变检测、圈速） \\
\texttt{saveSessionData/loadSessionData} & 用 ArduinoJson 把成绩存进 SPIFFS \\
\bottomrule
\end{tabularx}
\end{center}

\subsection{模块三：性能分析器 + 摩托车模式（performance_analyzer.cpp）}
\textbf{职责}：纯算法层。把 IMU + GPS 融合成 G 力，再衍生摩托专属指标。分两套数据
\texttt{fused\_data}（通用 G 力/功率）和 \texttt{moto\_data}（压弯/翘头/弯道）。
\begin{center}\small
\begin{tabularx}{\linewidth}{@{}lX@{}}
\toprule
\textbf{函数} & \textbf{作用} \\
\midrule
\texttt{calculateGForces()} & 动态信心度加权融合 GPS 差分加速度与 IMU；静止/动态两套逻辑 \\
\texttt{calculateLeanAngle()} & \textbf{互补滤波}算压弯倾角（陀螺仪积分 + 加速度计修正，\(\alpha\)=0.96） \\
\texttt{detectWheelieStoppie()} & 俯仰角 + 前向 G 力联合判断翘头/翘尾 \\
\texttt{analyzeCorner()} & 弯道进出检测 + 半径估算 $r = v^2/(g\cdot G_{lat})$ \\
\texttt{estimatePower()} & 滚阻+风阻+加速力模型估算功率(kW) \\
\bottomrule
\end{tabularx}
\end{center}
\begin{warn}{安装姿态硬编码}
\texttt{calculateLeanAngle()} 里 \texttt{Gyro\_Z\_Filtered} 被当作“横滚”、\texttt{Gyro\_X\_Filtered}
当作“偏航”，因为\textbf{传感器是垂直安装的}（X 轴朝下承重力）。坐标轴映射按这个安装姿态硬编码——
换安装方向就要改。
\end{warn}

\subsection{模块四：录音模块（audio_recorder.cpp）}
\textbf{职责}：定时器中断采集麦克风 ADC，存成 WAV 到 SPIFFS，骑行时自动录。
\textbf{数据结构}：\texttt{WAVHeader}（44 字节标准 WAV 头，先占位后回填 size）、
\texttt{AUDIO\_RECORDER\_DATA}（含录音状态机字段）。\textbf{设计思路}：8kHz/8bit/单声道——
刻意压低规格省 SPIFFS 空间。中断只读 ADC，写盘在主循环。

\subsection{模块五：SD 记录（sd_logger.cpp）}
\textbf{职责}：把每个数据点（18 字段）写成 CSV 行到 SD 卡，方便事后用 Excel/Python 分析。
\textbf{设计思路}：用 \textbf{HSPI}（而非默认 VSPI），MISO 改到 GPIO25 避开 GPIO12 启动冲突。
每 100 个点 flush 一次（平衡掉电安全与写入性能）。
\begin{warn}{半成品功能}
SD 初始化在 \texttt{setup()} 里被注释掉（“测试模式”），但 \texttt{loop()} 里仍调用
\texttt{checkAutoLogging} 和 \texttt{logDataPoint}。因 \texttt{initialized=false}，这些调用安全早退，不写盘。
\end{warn}

\subsection{模块六：Web 服务器（在 main.cpp 内）}
\textbf{职责}：WiFi 热点 + HTTP 服务，把数据 JSON 化推给手机网页。
\begin{center}\small
\begin{tabularx}{\linewidth}{@{}lX@{}}
\toprule
\textbf{路由} & \textbf{作用} \\
\midrule
\texttt{/} & 从 SPIFFS 返回 index.html（带 5 分钟缓存头） \\
\texttt{/data} & 主数据 JSON：温度/G力/角度/GPS/性能 \\
\texttt{/moto} & 摩托车模式 JSON：倾角/弯道/翘头/状态 \\
\texttt{/audio} & 录音状态 \\
\texttt{/audio/control} & 录音控制（start/stop/pause/list/...） \\
\texttt{/audio/download} & 下载 WAV 文件 \\
\bottomrule
\end{tabularx}
\end{center}
\textbf{设计思路}：用 \texttt{sprintf} 手拼 JSON（而非 ArduinoJson 序列化），固定 \texttt{char[]} 缓冲。
原因——\texttt{/data} 高频轮询，手拼比构建 JsonDocument 快、内存可控。代价是缓冲区大小要算准，否则溢出。

% ================================================================
\section{技术知识提取}
% ================================================================
\noindent 本部分只讲\textbf{项目里真的用到、且有“为什么”的}知识点。

\subsection{知识点 1：I²C（双路总线设计）}
MPU6050 与 OLED 都是 I²C 设备。项目用\textbf{两个独立的 I²C 控制器}：\texttt{TwoWire(0)} 给 MPU6050
（GPIO21/22），\texttt{TwoWire(1)} 给 OLED（GPIO18/19）。共享其实可行（地址 0x68/0x3C 不冲突），
分两路的好处是 OLED 刷新不会阻塞 IMU 的高频读取，代价是多占两个 GPIO。
\textbf{排错第一招}：\texttt{scanI2CDevices()} 分别扫两条总线。

\subsection{知识点 2：振动环境下的三级级联滤波（本项目灵魂）}
\textbf{解决}：发动机怠速 20\textasciitilde 50 Hz 振动叠加，G 力跳变、姿态抖 ±5°。
\begin{ascii}
原始ADC
  ① 硬件DLPF setFilterBandwidth(94Hz)   砍掉>94Hz，但保留快速动作响应
  ② 滑动平均FIR(窗口10@200Hz)            等效窗50ms，截止≈10Hz，无相位累积，延迟25ms
  ③ 一阶IIR低通(a=0.1)                   y=a*x+(1-a)*y_prev，截止≈3.5Hz，O(1)内存
  最终 *_Filtered (等效截止3~5Hz，振动衰减>12dB)
\end{ascii}
\textbf{关键权衡：}
\begin{itemize}
  \item DLPF 选 94Hz 而非 5Hz：陀螺仪要追摩托快速转向（可达 90°/s），截止太低动态响应变迟钝。硬件留宽，靠软件补。
  \item IIR 不用 FIR：MCU 上 IIR 只存一个上次输出值（O(1)），FIR 要存整个窗口。加速度测量不做频谱分析，对相位延迟不敏感。
  \item \(\alpha\)=0.1 怎么定：$f_c \approx \alpha f_s/(2\pi) \approx 0.1\times200/6.28 \approx 3.2$ Hz，再用串口波形微调。
\end{itemize}
\textbf{量化效果}：静止 G 力噪声 ±0.15g → ±0.01g；姿态抖动 ±5° → ±0.5°；代价 25\textasciitilde 50ms 延迟。

\subsection{知识点 3：互补滤波（Complementary Filter）}
\textbf{解决}：纯加速度计算姿态，动态时把“线加速度”误当“重力”，Pitch 虚假偏 5\textasciitilde 10°；
纯陀螺仪积分会漂移。\textbf{原理}：
\[ \theta = \alpha\,(\theta_{prev}+\omega\,dt) + (1-\alpha)\,\theta_{accel} \]
陀螺仪管高频（快、短期准），加速度计管低频（慢、长期准），互补。
\textbf{应用}：摩托倾角（\texttt{performance\_analyzer.cpp:524}）用 \(\alpha\)=0.96。注意\textbf{姿态主路径
（mpu6050.cpp）没用互补滤波}，是已知缺口——G 力优先级更高，姿态次要。
\textbf{误区}：\(\alpha\) 不是越大越好。\(\alpha\)=0.98 时间常数 $\tau=\alpha\cdot dt/(1-\alpha)\approx0.25$s，
这个“拉回速度”要和振动强度匹配着调。

\subsection{知识点 4：GPS-IMU 速度融合 + 动态信心度}
GPS 准但慢（10Hz、低速噪声大）；IMU 快（200Hz）但积分会漂。\textbf{两个融合点}：
\begin{enumerate}
  \item \textbf{速度融合}（\texttt{fuseIMUSpeed}，互补滤波）：200Hz 用加速度积分填 GPS 帧间空白，
        每个 GPS 帧（~100ms）到来时按 \texttt{GPS\_CORRECTION\_RATE=0.1} 拉回 10\% 漂移。
  \item \textbf{G 力融合}（\texttt{calculateGForces}，动态信心度）：按实时信号质量算
        \texttt{imu\_confidence = gps稳定性 × 信号质量}，动态分配 GPS 差分加速度与 IMU 的权重。
\end{enumerate}
\begin{center}\small
\begin{tabularx}{\linewidth}{@{}lXX@{}}
\toprule
 & \textbf{本项目(互补/自适应IIR)} & \textbf{卡尔曼(EKF)} \\
\midrule
计算量 & 极低，适合 ESP32 & 矩阵运算重 \\
调参 & 直观(\(\alpha\)、死区) & 要建模 Q/R 协方差 \\
选择 & [v] & 改进方向 \\
\bottomrule
\end{tabularx}
\end{center}
\textbf{误区}：用 GPS 速度\textbf{要用 Doppler 测速而非位置差分}。TinyGPS++ 的 \texttt{gps.speed.kmph()}
解析的 RMC 速度本身就是 Doppler 速度，精度比位置差分高 ~10 倍。项目用对了。

\subsection{知识点 5：GPS 零漂抑制（三层）}
GT-U7 静止时输出 0.5\textasciitilde 2 km/h 假速度。三层压制：
\begin{ascii}
层1：20点滑动平均 + 标准差分析   std_dev<0.2 且 avg<动态阈值 → 判定稳定
层2：自适应IIR (按速度段调a)     <0.5km/h→a=0.92强抑零漂；>4km/h→a=0.35快响应不耽误计时
层3：最终死区                    好信号(≥8星,HDOP<1.0)阈值0.8km/h，否则1.2km/h，以下归零
\end{ascii}
判定静止需持续 1.5 秒（15 次 ×100ms）。这是\textbf{精度 vs 响应}的反复权衡。

\subsection{知识点 6：ADC + 定时器中断采样（录音）}
\textbf{为什么用 ADC1 的 GPIO34}：ESP32 的 \textbf{ADC2 与 WiFi 冲突}（WiFi 开启时 ADC2 不可用）。
本项目全程开热点，麦克风必须接 ADC1。GPIO34 是纯输入引脚，最适合模拟输入。这是个容易踩、项目避开了的坑。
\textbf{采样}：硬件定时器 80 分频得 1MHz，125 tick=125µs=8kHz 触发，ISR 只 \texttt{adc1\_get\_raw} + 置标志。

\subsection{知识点 7：WAV 文件格式}
44 字节 RIFF 头 + PCM 裸数据。录音时先写占位头（size=0），结束时 \texttt{seek(0)} 回填真实
\texttt{dataSize}/\texttt{fileSize}。8bit 无符号 PCM，中心在 128。

\subsection{知识点 8：SPIFFS vs SD 卡}
\begin{center}\small
\begin{tabularx}{\linewidth}{@{}lXX@{}}
\toprule
 & \textbf{SPIFFS（板载Flash）} & \textbf{SD卡（外接）} \\
\midrule
容量 & ~1-3MB & GB级 \\
用途 & 网页文件 + 录音(小) & 行车CSV(大量) \\
接口 & 内部 & SPI(HSPI) \\
项目状态 & 在用 & 半成品(初始化被注释) \\
\bottomrule
\end{tabularx}
\end{center}

\subsection{知识点 9：PMTK / NMEA 协议}
GPS 用 NMEA 文本协议输出（\texttt{\$GPRMC}/\texttt{\$GPGGA}），用 PMTK 专有指令配置：
\texttt{\$PMTK220,100} 更新率 10Hz；\texttt{\$PMTK314,...} 只输出 RMC+GGA（防 9600 波特率溢出）；
\texttt{\$PMTK313,1} 启用 SBAS；\texttt{\$PMTK301,2} DGPS 用 WAAS。
\begin{warn}{文档与代码不一致}
\texttt{GT-U7\_ULTIMATE\_OPTIMIZATION.md} 说要把波特率提到 38400、用多星座，但 \texttt{initGPS()}
实际代码仍是 9600 + 10Hz。优化文档是“设想”，代码是“现状”。
\end{warn}

\subsection{知识点 10：HTTP / WiFi SoftAP}
ESP32 开热点（\texttt{WiFi.softAP}），固定 IP 192.168.4.1，跑 \texttt{WebServer}。网页 JS 轮询
\texttt{/data} 拿 JSON。\texttt{/data} 返回 \texttt{no-cache}（实时数据），\texttt{/}（HTML）返回
\texttt{max-age=300}（静态可缓存）——缓存策略区分对了。

% ================================================================
\section{设计决策分析}
% ================================================================
\begin{description}[leftmargin=0pt,style=nextline]
\item[决策 1：裸机分时调度，而非 FreeRTOS 任务]
  \texttt{loop()} 用 \texttt{millis()} 时间戳差比较，全程无 \texttt{delay()}。
  优点：简单、可控、无竞态、无切换开销；缺点：任一环节阻塞拖垮全局，没用上第二个核。
  数据流明确、无长阻塞操作，裸机分时足够且心智负担小。
\item[决策 2：全局数据结构通信，而非消息传递]
  零拷贝、无动态分配、实现极简；强耦合、不可重入（但单线程无此问题）。嵌入式裸机标准做法。
\item[决策 3：环形缓冲区做滑动平均]
  定长数组 + 取模索引 + \texttt{buffer\_full} 标志处理冷启动。O(1) 内存，无动态分配。
\item[决策 4：互补滤波而非卡尔曼]
  ESP32 算力和调参成本下，互补/自适应 IIR 性价比远高于 EKF。
\item[决策 5：双 I²C 总线]
  OLED 刷新不阻塞 IMU 读取；代价多占 2 个 GPIO。可单总线共享（地址不冲突）。
\item[决策 6：手拼 JSON 而非序列化库]
  热路径（\texttt{/data} 高频）选性能；冷路径（\texttt{saveSessionData}）用 ArduinoJson 选可维护——\textbf{两种都用，按场景分}。
\item[决策 7：SD 卡用 HSPI + MISO 改 GPIO25]
  GPIO12 是 strapping 引脚（启动拉高致 flash 电压判断错、启动失败），把 MISO 挪到 GPIO25 避开——ESP32 经典硬件坑。
\item[决策 8：录音 8kHz/8bit]
  SPIFFS 只有几 MB，录的是骑行环境音/排气声，不需高保真，容量优先（8KB/s，5 分钟约 2.4MB）。
\end{description}

% ================================================================
\section{调试与踩坑记录}
% ================================================================
\noindent 部分来自代码痕迹推断，部分已记录在 \texttt{INTERVIEW\_NOTES.md}。

\begin{center}\small
\begin{tabularx}{\linewidth}{@{}p{3.4cm}Xl@{}}
\toprule
\textbf{坑} & \textbf{原因 / 解决} & \textbf{状态} \\
\midrule
IMU 被发动机振动淹没 & DLPF 默认 260Hz 等于没滤；三级级联滤波，等效截止 3\textasciitilde 5Hz & [*]已解决 \\
GPS 静止零漂 & 多路径反射 + 几何误差；三层零漂抑制（均值+标准差/自适应IIR/死区） & [*]已解决 \\
IMU 偏置致积分漂移 & 出厂零偏 ±50\textasciitilde 150mg；手动硬编码 + GPS判静止时递推均值在线校准 & [*]部分解决 \\
GPS 位置跳变 & 遮挡后重定位瞬跳数百米；Haversine 算帧间距离，>1000m 拒绝 & [*]已解决 \\
GPIO12 启动冲突(SD) & strapping 引脚；MISO 挪到 GPIO25，用 HSPI & [*]已规避 \\
ADC2 与 WiFi 冲突(麦克风) & WiFi 开启时 ADC2 读不到；麦克风接 ADC1 的 GPIO34 & [*]已规避 \\
9600 下 10Hz NMEA 溢出 & 全量语句溢出丢帧；PMTK314 只留 RMC+GGA & [*]已规避 \\
文档与代码脱节 & README/优化文档 与代码 多处对不上（见下） & [!] 待修 \\
SD 记录是半成品 & initSDCard() 被注释，loop() 仍调（initialized=false 安全早退） & [!] 待完成 \\
\bottomrule
\end{tabularx}
\end{center}

\subsubsection{脱节细节（坑 8 展开）}
\begin{itemize}
  \item README 写“共享 I²C 总线 GPIO21/19”，实际代码是\textbf{双路} GPIO21/22 + GPIO18/19。
  \item README 的 \texttt{lib\_deps} 写 \texttt{electroniccats/MPU6050}，实际 \texttt{platformio.ini} 用 \texttt{adafruit/Adafruit MPU6050}。
  \item \texttt{GT-U7\_ULTIMATE} 说 38400 波特率/多星座，代码实际 9600/10Hz。
\end{itemize}

\subsubsection{潜在风险（代码审查角度，未必触发）}
\begin{itemize}
  \item \textbf{JSON 缓冲区}：\texttt{/data} 用 \texttt{char[750]}、\texttt{/moto} 用 \texttt{char[600]} 手拼，
        字段多、\texttt{\%.8f} 经纬度长，\textbf{接近上限，加字段易溢出}。建议加 \texttt{snprintf} 长度保护。
  \item \textbf{静止校准 static 局部变量}：\texttt{calculateGForces} 里 \texttt{gravity\_bias\_x} 等用 \texttt{static}
        且只累加不重置，长时间运行逻辑可能僵化（满 50 样本后固定）。
\end{itemize}

% ================================================================
\section{工程实践总结}
% ================================================================
\subsubsection{做得好的地方}
\begin{itemize}
  \item \textbf{代码组织}：按外设/算法清晰分模块，头文件声明 + cpp 实现规整。
  \item \textbf{错误处理}：初始化大量带降级（MPU6050 试 0x68 失败试 0x69；OLED 失败仍继续；GPS 信号丢失自动热重启）。\textbf{失败不死机、打印原因、降级运行}的思路很成熟。
  \item \textbf{可观测性}：串口三种打印模式可运行时切换；I²C 扫描、原始 NMEA 打印——调试钩子完备。
  \item \textbf{构建系统}：PlatformIO 依赖锁版本，\texttt{monitor\_filters = esp32\_exception\_decoder} 自动解崩溃栈。
  \item \textbf{配置管理}：引脚、阈值、滤波系数都用 \texttt{\#define}/\texttt{const} 集中在文件头。
  \item \textbf{场景化分层}：热路径手拼 JSON 求性能，冷路径用 ArduinoJson 求可维护——分得清。
\end{itemize}
\subsubsection{可以改进的地方}
\begin{itemize}
  \item \textbf{文档同步}：README/优化文档已严重脱节（坑 8）。
  \item \textbf{魔法数字}：\texttt{-0.47/+0.5/-0.48} 校准偏置散落多文件，应抽成统一配置或自动校准。
  \item \textbf{缓冲区安全}：手拼 JSON 用 \texttt{sprintf} 无长度保护，应换 \texttt{snprintf}。
  \item \textbf{SD 功能收尾}：要么启用要么删，别留半成品调用。
  \item \textbf{测试策略}：目前\textbf{无任何自动化测试}（\texttt{test/} 只有空 README）。IMU 滤波、Haversine、零漂判定这些纯函数很适合 host 端单元测试。
  \item \textbf{持续集成}：无 CI，可加 \texttt{pio ci} 在 push 时编译检查。
\end{itemize}

% ================================================================
\section{项目复盘}
% ================================================================
\subsubsection{最大的亮点}
把“算法”当成项目主角，而非“点亮外设”。\textbf{三级滤波 + GPS-IMU 融合 + 动态信心度}这套组合，
把 30 元 GPS 的静止零漂从 ±2.5 压到 ±0.8 km/h（与专业设备差距从 20 倍缩到 4 倍，成本差 67 倍）。

\subsubsection{最大的难点}
\textbf{在脏数据上做可信的实时融合}。难点不在写公式，而在：什么时候信谁、信多少（动态信心度）、
怎么不让滤波延迟拖垮测试计时（自适应 \(\alpha\)）。都是反复调参 + 串口看波形磨出来的。

\subsubsection{学到的知识}
滤波器选型的工程权衡（FIR vs IIR、硬件 vs 软件、截止频率 vs 响应延迟）；传感器融合实战；
ESP32 硬件陷阱（strapping 引脚、ADC2/WiFi 冲突）；裸机分时调度模式；GPS NMEA/PMTK、Doppler 测速。

\subsubsection{如果重新设计会怎么做}
\begin{itemize}
  \item 传感器融合上 \textbf{1D 卡尔曼}（只融合纵向速度，状态量少，ESP32 扛得住）。
  \item 校准全自动化，删硬编码偏置，统一静止在线校准 + 温度补偿。
  \item 姿态主路径加互补滤波，解决加速时 Pitch 虚假偏移。
  \item 抽配置层（所有引脚/阈值/系数集中到 \texttt{config.h}）。
  \item 加 host 端单元测试；用 FreeRTOS 双核物理隔离 Web/录音 与 传感器融合。
\end{itemize}

\subsubsection{下一步可以扩展什么}
升级 GPS 到 NEO-M9N（\texttt{GPS\_UPGRADE\_ANALYSIS.md} 已选型，150 元换 5-8 倍精度）；真正启用 SD 记录 +
Python 分析脚本；加磁力计（QMC5883L）补全 Yaw；大倾角重力分量补偿 $a_{\mathrm{true}}=a_{\mathrm{meas}}-g\sin(\mathrm{Roll})$；
网页端轨迹回放 / 成绩对比。

% ================================================================
\section{个人知识地图}
% ================================================================
\noindent 由本工作区内容沉淀出的长期知识体系。带 [v] 的是本项目实际用过的。
\begin{ascii}
Embedded Knowledge Map
│
├── MCU 基础 (ESP32)
│   ├── GPIO [v]           (BOOT按钮输入、CS控制)
│   ├── ADC [v]            (ADC1_CH6麦克风；ADC2/WiFi冲突坑)
│   ├── Timer + 中断 [v]    (8kHz采样定时器、IRAM_ATTR、ISR最小化)
│   ├── Strapping引脚 [v]   (GPIO12启动冲突坑)
│   └── 双核 (未用，改进方向)
│
├── 通信协议
│   ├── I²C [v]            (双路总线、设备扫描、400kHz、地址0x68/0x3C)
│   ├── UART [v]           (GPS 9600、HardwareSerial2)
│   ├── SPI [v]            (SD卡 HSPI、避VSPI冲突)
│   ├── HTTP/WiFi [v]      (SoftAP热点、WebServer、JSON、缓存策略)
│   └── NMEA/PMTK [v]      (GPS文本协议、Doppler测速、配置指令)
│
├── 传感器与信号处理 [*]本项目核心
│   ├── IMU (MPU6050) [v]  (量程/DLPF配置、6轴)
│   ├── 滤波器 [v]
│   │   ├── 硬件DLPF [v]      ├── 滑动平均(FIR) [v]
│   │   ├── 一阶IIR/EMA [v]   └── 卡尔曼 (未用，改进方向)
│   ├── 传感器融合 [v]
│   │   ├── 互补滤波 [v](倾角,a=0.96)  ├── 动态信心度加权 [v]  └── 阈值穿越插值 [v]
│   ├── 姿态解算 [v]       (atan2算Roll/Pitch；Yaw缺磁力计)
│   └── 偏置校准 [v]       (递推均值O(1)在线校准)
│
├── GPS / 定位
│   ├── NMEA解析 [v](TinyGPS++)   ├── 零漂抑制 [v](三层)
│   ├── Haversine距离 [v]          ├── HDOP/卫星质量 [v]   └── SBAS/DGPS增强 [v]
│
├── 数据存储
│   ├── SPIFFS [v](网页/录音/会话JSON)  ├── SD卡 CSV [!](半成品)
│   ├── WAV格式 [v]                     └── ArduinoJson [v]
│
├── 系统架构 / 软件工程
│   ├── 裸机分时调度 [v](millis时间片,无delay)   ├── 全局单例通信 [v]
│   ├── 环形缓冲区 [v]   ├── 状态机 [v](录音/测试/SD)
│   ├── 错误降级处理 [v] └── RTOS任务/队列 (未用,改进方向)
│
├── 车辆动力学 (应用领域)
│   ├── G力(纵向/横向/垂直) [v]   ├── 压弯倾角 [v]   ├── 弯道半径 r=v²/(gG) [v]
│   ├── 翘头/翘尾检测 [v]          ├── 加速/制动测试 [v]  └── 功率估算(阻力模型) [v]
│
└── 机械 / 工具链
    ├── PlatformIO [v]   ├── OpenSCAD参数化建模 [v](cad/外壳)   └── Gerber/PCB [v]
\end{ascii}

\vspace{8pt}
\begin{center}\small\color{gray}\itshape
本知识库由对工作区源码、文档、配置的递归分析生成。凡代码无法确定之处均已标注；未发现编造的功能或设计。\\
若代码变更，请同步更新本文档——过时的笔记比没有笔记更危险。
\end{center}

\end{document}
